Application workflow and the actual archival learning target. (a) Three slots from one archived ISO~834 case are projected schematically; matching colours identify the slots and graph nodes. The magnified opening retains the source dimensions. (b) The graph and prescribed-exposure branches feed a scalar decoder. The seven comparison variants use the same experimental protocol; graph transfer tests retain the trained weights when evaluating unseen graph sizes or alternative radius/symmetric-kNN views. The prediction symbol indicates the task, not a plotted model result. (c) The same case supplies 571 native scalar observations. Their running minimum is linearly interpolated at 61 times from 0 to 3600~s and stored in float32, exactly following the frozen archival-target preparation. The enlarged interval makes the scalar-envelope operation visible. All curves in (c) are source observations or deterministic target preprocessing; no predicted curve is shown. This 997-case benchmark includes 86 restored constant responses and uses an archival scalar reference rather than a reconstructed pointwise thermal-capacity label.
